% Standalone TikZ figure: BOHB (Bayesian Optimization with Hyperband) Algorithm
% For Chapter 3: Reinforcement Learning for Parameter Optimization
\documentclass[tikz,border=10pt]{standalone}
\usepackage{tikz}
\usepackage{amsmath}
\usetikzlibrary{arrows.meta, positioning, shapes.geometric, fit, backgrounds, calc, patterns, decorations.pathreplacing}

\begin{document}
\begin{tikzpicture}[
    node distance=0.8cm and 1.2cm,
    every node/.style={font=\small},
    configbox/.style={rectangle, draw, minimum width=0.6cm, minimum height=0.6cm, rounded corners=2pt},
    arrow/.style={-{Stealth[length=2mm]}, thick},
    dashedarrow/.style={-{Stealth[length=2mm]}, thick, dashed},
    bracketlabel/.style={font=\footnotesize\bfseries},
    label/.style={font=\footnotesize},
    title/.style={font=\bfseries}
]

% Title
\node[title, font=\large] at (5, 8) {BOHB: Multi-Fidelity Successive Halving};

% ============ BRACKET VISUALIZATION ============

% Bracket 0 (most aggressive - many configs, low budget start)
\node[bracketlabel] at (-1.5, 6.5) {Bracket $s=2$};
\node[label, text=gray] at (-1.5, 6) {(Aggressive)};

% Initial configs for bracket 0 (9 configs)
\foreach \i in {1,...,9} {
    \node[configbox, fill=blue!30] (b0r0c\i) at ({0.8*\i}, 6) {};
}
\node[label] at (10, 6) {$r_0 = R/9$};

% After first halving (3 configs)
\foreach \i in {1,...,3} {
    \node[configbox, fill=blue!50] (b0r1c\i) at ({0.8*\i + 2.4}, 4.8) {};
}
\node[label] at (10, 4.8) {$r_1 = R/3$};

% After second halving (1 config)
\node[configbox, fill=blue!80] (b0r2c1) at (4, 3.6) {};
\node[label] at (10, 3.6) {$r_2 = R$};

% Arrows showing elimination
\draw[arrow, red!60] (b0r0c1.south) -- (b0r1c1.north);
\draw[arrow, red!60] (b0r0c4.south) -- (b0r1c2.north);
\draw[arrow, red!60] (b0r0c7.south) -- (b0r1c3.north);
\draw[arrow, red!60] (b0r1c2.south) -- (b0r2c1.north);

% X marks for eliminated configs
\foreach \i in {2,3,5,6,8,9} {
    \node[text=red!70, font=\small] at (b0r0c\i.center) {$\times$};
}
\node[text=red!70, font=\small] at (b0r1c1.center) {$\times$};
\node[text=red!70, font=\small] at (b0r1c3.center) {$\times$};

% ============ TPE SAMPLING SECTION ============
\node[title] at (-3, 3) {TPE Sampling};

% Good/bad distribution visualization
\node[draw, rounded corners, fill=green!10, minimum width=2.5cm, minimum height=1.5cm] (good) at (-3, 1) {};
\node[label, text=green!60!black] at (-3, 1.5) {$\ell(\mathbf{x})$ (good)};
\draw[green!60!black, thick, domain=-1:1, samples=30] plot ({-3 + 0.8*\x}, {0.5 + 0.8*exp(-4*\x*\x)});

\node[draw, rounded corners, fill=red!10, minimum width=2.5cm, minimum height=1.5cm] (bad) at (-3, -1.2) {};
\node[label, text=red!60!black] at (-3, -0.7) {$g(\mathbf{x})$ (bad)};
\draw[red!60!black, thick, domain=-1:1, samples=30] plot ({-3 + 0.8*\x}, {-1.7 + 0.5*(1 + 0.3*sin(720*\x))});

% Acquisition function
\node[label] at (-3, -2.8) {Maximize: $\displaystyle\frac{\ell(\mathbf{x})}{g(\mathbf{x})}$};

% Arrow from TPE to configs
\draw[dashedarrow, blue!70] (-1.5, 1) -- (0.8, 6) node[pos=0.5, above, sloped, font=\footnotesize] {sample configs};

% ============ HYPERBAND BRACKET STRUCTURE ============
\node[title] at (7, 2) {Bracket Structure};

% Table-like structure for brackets
\node[draw, rounded corners, fill=gray!10, minimum width=5cm, minimum height=3.5cm] at (7, 0) {};

% Headers
\node[label, font=\bfseries] at (5.3, 1.3) {Bracket};
\node[label, font=\bfseries] at (6.5, 1.3) {$n$ (configs)};
\node[label, font=\bfseries] at (7.7, 1.3) {$r_0$ (budget)};
\node[label, font=\bfseries] at (8.8, 1.3) {Strategy};

% Bracket s=2
\node[label] at (5.3, 0.6) {$s=2$};
\node[label] at (6.5, 0.6) {9};
\node[label] at (7.7, 0.6) {$R/9$};
\node[label, text=blue!70] at (8.8, 0.6) {Explore};

% Bracket s=1
\node[label] at (5.3, 0) {$s=1$};
\node[label] at (6.5, 0) {3};
\node[label] at (7.7, 0) {$R/3$};
\node[label, text=purple!70] at (8.8, 0) {Balanced};

% Bracket s=0
\node[label] at (5.3, -0.6) {$s=0$};
\node[label] at (6.5, -0.6) {1};
\node[label] at (7.7, -0.6) {$R$};
\node[label, text=red!70] at (8.8, -0.6) {Exploit};

% Lines
\draw[gray] (4.5, 1) -- (9.3, 1);
\draw[gray] (4.5, 0.3) -- (9.3, 0.3);
\draw[gray] (4.5, -0.3) -- (9.3, -0.3);

% ============ FEEDBACK LOOP ============
\node[draw, rounded corners, fill=orange!15, minimum width=2cm, minimum height=1cm] (history) at (2, -1.5) {History $\mathcal{H}$\\$\{(\mathbf{x}, y, r)\}$};

% Arrows showing feedback
\draw[arrow, orange!70] (b0r2c1.south) -- ++(0, -1) -| (history.north);
\draw[dashedarrow, orange!70] (history.west) -- ++(-1, 0) |- (good.east);

% Labels for process
\node[label, text=gray, align=center] at (2, -3.2) {Update TPE models\\after evaluations};

% ============ LEGEND ============
\node[anchor=north west] at (-4.5, -3.5) {
    \begin{tikzpicture}[every node/.style={font=\footnotesize}]
        \node[configbox, fill=blue!30] at (0,0) {};
        \node[right] at (0.4, 0) {Low-fidelity evaluation};
        \node[configbox, fill=blue!80] at (0,-0.5) {};
        \node[right] at (0.4, -0.5) {Full-fidelity evaluation};
        \node[text=red!70] at (0.1, -1) {$\times$};
        \node[right] at (0.4, -1) {Eliminated (bottom performers)};
    \end{tikzpicture}
};

% Reduction factor annotation
\draw[decorate, decoration={brace, amplitude=6pt, raise=2pt, mirror}] (b0r0c1.south west) -- (b0r0c9.south east)
    node[midway, below=10pt, font=\footnotesize] {Reduce by $\eta=3$ each round};

\end{tikzpicture}
\end{document}
